\section{Introduction}

Offline goal-conditioned reinforcement learning (GCRL) provides a simple and fully self-supervised recipe for learning reusable behavior from unlabeled
trajectories: relabel states in the dataset as goals, learn a goal-conditioned value or representation, and extract a policy that reaches arbitrary dataset
states. Recent benchmarks such as OGBench have made this setting substantially more systematic by evaluating a diverse set of algorithms across locomotion,
manipulation, stochastic, visual, and long-horizon tasks. However, current benchmark practice still largely reports final performance after tuning. This
leaves a basic question unanswered: \emph{How easily can a policy be extracted from a learned value function?}

This question is especially important in offline RL because policy extraction is not a minor implementation detail.
Most methods regularize the learned policy toward dataset actions, and the strength of this regularization is often one of the most important hyperparameters.
In advantage-weighted regression (AWR), the actor is trained by weighted behavioral cloning and the weights are determined by the advantages.
Thus, policy extraction depends not only on whether the value model
is accurate, but on whether its advantages are \emph{well-conditioned}: whether they induce a stable and useful weighting of dataset actions under exponentiation.

The difficulty lies in the fact that different offline GCRL algorithms assign fundamentally different semantics to this advantage: bootstrapped temporal-difference margins (GCIQL), local value progress (GCIVL), quasimetric distance reduction (QRL), or contrastive reachability ratios (CRL).
Because these semantics vary so widely, evaluating them solely on success metrics can mischaracterize their utility.
What ultimately matters for the AWR extractor is not a universal notion of distance, but whether the learned signal reliably separates useful dataset actions from irrelevant ones.

\begin{figure}[t]
  \centering
  \begin{minipage}[c]{0.66\textwidth}
    \centering
    \begin{tabular}{@{}cc@{}}
      \includegraphics[width=0.48\linewidth]{figures/awr-mobility/antmaze-medium/90/mobility-gciql-antmaze-medium-navigate-v0.png}
      &
      \includegraphics[width=0.48\linewidth]{figures/awr-mobility/antmaze-medium/90/mobility-gcivl-antmaze-medium-navigate-v0.png}
      \\
      \footnotesize (a) GCIQL & \footnotesize (b) GCIVL \\ [4pt]
      \includegraphics[width=0.48\linewidth]{figures/awr-mobility/antmaze-medium/90/mobility-crl-antmaze-medium-navigate-v0.png}
      &
      \includegraphics[width=0.48\linewidth]{figures/awr-mobility/antmaze-medium/90/mobility-qrl-antmaze-medium-navigate-v0.png}
      \\
      \footnotesize (c) CRL & \footnotesize (d) QRL
    \end{tabular}
  \end{minipage}\hfill
  \begin{minipage}[c]{0.32\textwidth}
    \captionsetup{
      type=figure,
      justification=raggedright,
      singlelinecheck=false,
      font=small
    }
    \caption{\textbf{Trainability landscapes in the learning-rate ($\eta$) -- AWR-temperature ($\alpha$) plane.}
    We visualize the top-10\% regions to compare how different objectives interact with the policy extractor.
    \textbf{Progress-based methods} (GCIVL, QRL) exhibit broad, stable winner sets across a wide range of $\alpha$, indicating high extractability.
    \textbf{CRL} shows a concentrated but stable region, suggesting its reachability signal supports selective updates.
    \textbf{GCIQL} (bootstrapped $Q-V$) displays sharper, more mobile landscapes that are highly sensitive to both hyperparameters, reflecting the difficulty of extracting from bootstrapped margins.
    }
    \label{fig:mobility:phases_vs_quality-antmaze}
  \end{minipage}
\end{figure}

We therefore study offline GCRL through \emph{trainability landscapes}.
Following the hyperparameter-landscape methodology of \citet{mohan-automlconf23a}, we view training as a phase-indexed response surface \(f_t:\Lambda\to\mathbb{R}\), where each hyperparameter configuration \(\lambda\in\Lambda\) maps to downstream evaluation performance at phase \(t\).
Unlike prior online RL landscape studies, we work in a fixed offline dataset setting, which removes exploration-induced distribution shift.
Unlike standard algorithm benchmarks, we do not only ask which method achieves the highest success; we ask which methods expose broad, stable near-optimal regions in the hyperparameter space, and what actor-facing signals make such regions possible.
Our primary analytical lens is the shared $(\eta,\alpha)$ landscape, capturing the central tension of offline extraction: the learning rate ($\eta$) controls the shape of the learned value signal, while the AWR temperature ($\alpha$) controls how aggressively that signal is converted into a policy.
By holding the extractor fixed, we ask not just which method peaks highest, but which methods expose broad, stable near-optimal regions.

% Our main experimental lens is the shared \((\eta,\alpha)\) landscape, where \(\eta\) is the learning rate and \(\alpha\) is the AWR temperature. This pair
% captures the central interaction in offline GCRL extraction: the learning rate controls the value or representation signal that is learned, while \(\alpha\)
% controls how aggressively that signal is converted into a policy. By holding the extractor fixed across algorithms, we can compare value-learning schemes through
% the geometry they induce for the actor. We find that these geometries differ substantially: some methods produce broad positive support, others produce calibrated but sparse winners, and others produce sparse or miscalibrated advantages whose extraction is highly temperature-sensitive.
Our analysis reveals three qualitative regimes.
First, progress-based methods such as GCIVL and QRL often produce broad winner sets: many dataset transitions
receive positive advantages, making AWR extraction comparatively diffuse.
Second, CRL often produces fewer positive advantages, but the successful configurations use highly selective AWR updates: lower effective sample size correlates with better return, suggesting that its contrastive reachability signal supports stable hard selection of dataset actions.
Third, GCIQL exhibits sharper and more mobile landscapes, including distinct regimes in the learning-rate--temperature plane.
We hypothesize this to be due to the difficulty of extracting from a bootstrapped \(Q-V\) margin whose positive mass is not always predictive of return.

These findings suggest that final return alone misses an important benchmark axis: \emph{extractability}.
A value-learning scheme may be accurate in its own semantics, yet hard to use with a behavioral policy extractor if its
actor-facing signal is poorly calibrated, overly sparse, or highly sensitive to temperature.
Conversely, a method such as CRL may be stable not because it estimates geometric distance, but because it learns a behavior-supported
reachability ranking that AWR can safely concentrate on.
To bridge this gap, we introduce diagnostics connecting response-surface geometry to the learned objective.

\paragraph{Contributions.}
Our contributions are:
\begin{inparaenum}[(i)]
    \item We introduce trainability landscapes for offline GCRL under a shared policy extractor, focusing on the learning-rate--AWR-temperature plane as
    the central interaction between value learning and behavioral policy extraction.
    \item We show that different value-learning schemes induce qualitatively different extraction geometries, yielding broad, calibrated, or brittle hyperparameter landscapes even on fixed offline datasets.
    \item We identify advantage extractability as a mechanism behind this behavior: stable landscapes arise when AWR sees a reliable winner set, while sparse or miscalibrated advantages induce stronger temperature sensitivity.
    \item We propose lightweight diagnostics, including AWR winner-set geometry and future-vs-random advantage separation, that help explain when learned advantages are reliable for policy extraction.
\end{inparaenum}
